\documentclass[10pt]{article}

\usepackage[T1]{fontenc}
\usepackage[utf8]{inputenc}
\usepackage{lmodern}
\usepackage{microtype}
\usepackage[margin=2.2cm]{geometry}
\usepackage{amsmath,amssymb}
\usepackage{booktabs}
\usepackage{longtable}
\usepackage{tabularx}
\usepackage{array}
\usepackage{multirow}
\usepackage{enumitem}
\usepackage{xcolor}
\usepackage[hidelinks]{hyperref}
\usepackage{caption}

\setlist{itemsep=0.15em,topsep=0.3em}
\setlength{\parskip}{0.35em}
\setlength{\parindent}{0pt}
\captionsetup{font=small,labelfont=bf}
\newcolumntype{L}[1]{>{\raggedright\arraybackslash}p{#1}}
\newcommand{\code}[1]{\texttt{#1}}
\newcommand{\status}[1]{\textcolor{black!55}{\textit{[#1]}}}

\title{\textbf{BrainModel: a modality-agnostic 3D foundation model for brain CT and MRI}\\[0.3em]
\large Setup, data, preprocessing, methods and planned experiments}
\author{Working document}
\date{October 2026}

\begin{document}
\maketitle

\begin{abstract}
We are building a single self-supervised 3D encoder for brain CT and MRI whose frozen embeddings should serve many
downstream tasks: hemorrhage and stroke detection, tumor and lesion segmentation, dementia staging, brain age and
more. This document fixes the setting, the data (about 30 public sources, roughly 110{,}000 volumes), the shared
preprocessing that places every scan in one anatomical frame, the model and training objective, the downstream
evaluation protocol, and the experiments planned to finalize the model and then to post-train it with task
conditioning.
\end{abstract}

\tableofcontents
\bigskip

%=====================================================================================================================
\section{Motivation}

Brain imaging AI is fragmented: most models are trained for one task, one modality and often one scanner. Yet CT and
the many MRI contrasts image the same anatomy, and large amounts of unlabeled scans are available while labels are
scarce, heterogeneous and expensive. A shared encoder pretrained without labels on both modalities could (i) give
strong frozen features for new tasks with little labeled data, (ii) transfer knowledge across modalities, and (iii)
provide one representation that many clinical questions can be read from. The design goal is therefore generality:
an embedding that is useful across many tasks, not the best model on any single one.

%=====================================================================================================================
\section{Setting}

\paragraph{Self-supervised pretraining.} A 3D Vision Transformer is trained on unlabeled CT and MRI volumes with a
joint-embedding objective (LeJEPA~\cite{lejepa}): different crops and intensity renderings of the same volume must map
to nearby embeddings, while a distribution penalty (SIGReg) keeps the embedding space isotropic and prevents collapse.
Agreement is enforced both for a pooled (volume-level) embedding and token by token, and a masked-prediction term
keeps local detail. Labels are never used, and downstream evaluation families are held out of pretraining entirely.

\paragraph{Frozen evaluation.} The pretrained encoder is frozen. Every volume is embedded once (full-volume token
maps), and light heads are fit on the embeddings for each labeled task: linear/ridge readouts, a small 3D
convolutional probe, and a token decoder for segmentation. Token-level probes measure what the representation
encodes without any task head (lesion detectability, scanner/dataset fingerprint, anatomical position). A
randomly initialized encoder of the same architecture is evaluated identically and serves as the floor.

\paragraph{Task-conditioned post-training.} After the self-supervised model is finalized (size, hyperparameters,
training length, data selection and preprocessing), it is post-trained jointly on many labeled tasks, conditioned on
a learned task vector that identifies the task (e.g.\ age regression, hemorrhage classification). The post-trained
model is compared with the frozen model plus heads, on the tasks it was trained on and, crucially, on tasks held out
of post-training.

\paragraph{Experiment records.} Every evaluation writes a self-contained record (model and checkpoint provenance,
data version, protocol, validation and test metrics with confidence intervals, environment, artifacts) that is
ingested into one SQLite database, so models, heads, seeds and data versions are compared on identical footing.

%=====================================================================================================================
\section{Data}

\subsection{Sources}

Table~\ref{tab:data} lists every dataset. \emph{Role} is \textbf{SSL} (unlabeled pretraining; the dataset may also
carry downstream labels) or \textbf{held out} (never seen in pretraining; evaluation only). Counts are volumes (one
volume = one 3D scan of one sequence) and subjects in the current store; they are updated by the final index build.

\begingroup
\footnotesize
\setlength{\tabcolsep}{3.5pt}
\begin{longtable}{L{2.35cm} L{5.5cm} L{1.95cm} r r L{1.25cm} L{2.6cm}}
\caption{Datasets. Mod.: modality and sequences in the store. Vol./Subj.: volumes and subjects.}\label{tab:data}\\
\toprule
Dataset & Description & Mod. & Vol. & Subj. & Role & Tasks \\
\midrule
\endfirsthead
\toprule
Dataset & Description & Mod. & Vol. & Subj. & Role & Tasks \\
\midrule
\endhead
\midrule \multicolumn{7}{r}{\textit{continued}} \\
\endfoot
\bottomrule
\endlastfoot
\multicolumn{7}{l}{\textbf{CT}}\\
RSNA ICH~\cite{rsna} & Head CT studies from the RSNA 2019 Intracranial Hemorrhage Detection challenge, labeled per slice for any hemorrhage and five subtypes. & CT & 25{,}137 & 25{,}137 & SSL & \code{rsna\_ich} \\
CQ500~\cite{cq500} & Non-contrast head CTs from Indian hospitals, read by three radiologists for hemorrhage subtypes, calvarial fracture, mass effect and midline shift. & CT & 422 & 422 & held out & \code{cq500} \\
Seg-CQ500 & A CQ500 subset with voxel-level intracranial hemorrhage masks. & CT & 51 & 51 & held out & \code{seg\_cq500}, \code{cq500} \\
INSTANCE 2022 & Non-contrast head CTs with intracranial hemorrhage masks from the MICCAI 2022 INSTANCE challenge. & CT & 100 & 100 & held out & \code{instance} \\
AISD~\cite{aisd} & Non-contrast CTs of acute ischemic stroke with infarct masks drawn with reference to follow-up diffusion MRI. & CT & 398 & 398 & held out & \code{aisd} \\
ISLES 2024 & Acute stroke admission CT (non-contrast and CT angiography) with infarct masks from follow-up MRI. & NCCT, CTA & 298 & 149 & held out & \code{isles24} \\
\midrule
\multicolumn{7}{l}{\textbf{MRI: tumors}}\\
BraTS GLI (2023, 2024)~\cite{brats21} & Adult glioma, pre-operative (2023) and post-treatment (2024), four co-registered skull-stripped sequences with tumor sub-region labels. & T1, T1-Gd, T2, FLAIR & 11{,}488 & 1{,}830 & SSL & \code{brats2023\_gli}, \code{brats2024\_gli} \\
BraTS MEN 2023~\cite{bratsmen} & Pre-operative meningioma, four skull-stripped sequences with tumor labels. & T1, T1-Gd, T2, FLAIR & 4{,}000 & 944 & SSL & \code{brats2023\_men} \\
BraTS PED 2023~\cite{bratsped} & Pediatric brain tumors (mostly high-grade gliomas), four skull-stripped sequences with tumor labels. & T1, T1-Gd, T2, FLAIR & 576 & 144 & SSL & \code{brats2023\_ped} \\
BraTS-GoAT & Mixed tumor types from the BraTS generalizability-across-tumors cohort, four skull-stripped sequences (via the FOMO corpus). & T1, T1-Gd, T2, FLAIR & 10{,}444 & 2{,}611 & SSL & -- \\
MSD Brain Tumour~\cite{msd} & Multi-parametric glioma MRIs from the Medical Segmentation Decathlon (Task 01); the T1 channel is used. & T1 & 484 & 484 & held out & \code{msd} \\
BrainMetShare~\cite{brainmetshare} & Skull-stripped multi-sequence MRI of patients with brain metastases (Stanford). & T1, T1-Gd, BRAVO, FLAIR & 624 & 156 & held out & reserved \\
FOMO Yale brain mets & Longitudinal, skull-stripped multi-sequence MRI of patients with brain metastases. & T1, T1-Gd, T2, FLAIR, PD & 33{,}681 & 1{,}430 & SSL & -- \\
\midrule
\multicolumn{7}{l}{\textbf{MRI: stroke and white-matter disease}}\\
ATLAS~\cite{atlas} & T1 MRI of chronic stroke with manually traced lesion masks. & T1 & 1{,}453 & 1{,}453 & held out & \code{atlas} \\
ISLES 2022~\cite{isles22} & Acute and sub-acute ischemic stroke MRI with lesion masks in diffusion space. & DWI, ADC, FLAIR & 750 & 250 & held out & \code{isles22} \\
WMH 2017~\cite{wmh} & FLAIR and T1 MRI from several scanners with white-matter hyperintensity masks (MICCAI 2017 challenge). & FLAIR, T1 & 340 & 170 & held out & \code{wmh} \\
\midrule
\multicolumn{7}{l}{\textbf{MRI: aging, dementia, development and healthy cohorts}}\\
OASIS-1~\cite{oasis1} & Cross-sectional T1 MRI of adults aged 18--96, including very mild to mild Alzheimer's disease, with CDR. & T1 & 436 & 416 & SSL & \code{oasis1\_age}, \code{oasis1\_cdr} \\
OASIS-2~\cite{oasis2} & Longitudinal T1 MRI of nondemented and demented older adults (60--96) with CDR at each visit. & T1 & 373 & 150 & SSL & \code{oasis2\_cdr}, \code{oasis2\_conversion} \\
IXI & T1 and T2 MRI of about 600 healthy adults from three London hospitals, with demographics. & T1, T2 & 1{,}159 & 582 & SSL & \code{ixi\_age} \\
ABIDE I~\cite{abide} & T1 MRI of individuals with autism spectrum disorder and typically developing controls from 17 sites. & T1 & 1{,}102 & 1{,}102 & SSL & \code{abide\_asd} \\
AOMIC PIOP1~\cite{aomic} & T1 MRI of young adults (Amsterdam Open MRI Collection, population imaging of psychology 1). & T1 & 216 & 216 & SSL & \code{piop1\_age} \\
AOMIC PIOP2~\cite{aomic} & As PIOP1, a second cohort. & T1 & 226 & 226 & SSL & \code{piop2\_age} \\
AOMIC ID1000~\cite{aomic} & T1 MRI of about 1{,}000 young adults with broad demographic and psychometric data. & T1 & 928 & 928 & SSL & \code{id1000\_age} \\
Pixar~\cite{pixar} & T1 MRI of children aged 3--12 and adults from a movie-viewing fMRI study (OpenNeuro ds000228). & T1 & 155 & 155 & SSL & \code{pixar\_age} \\
\midrule
\multicolumn{7}{l}{\textbf{MRI: FOMO pretraining corpus (structural sequences of selected families)}}\\
HBN~\cite{hbn} & Healthy Brain Network: children and adolescents, a transdiagnostic pediatric mental-health cohort. & T1, T2, FLAIR, PD & 8{,}481 & 3{,}843 & SSL & -- \\
NKI-Rockland~\cite{nki} & Community lifespan sample from the Nathan Kline Institute. & T1, T2, FLAIR & 5{,}951 & 1{,}326 & SSL & -- \\
CoRR~\cite{corr} & Consortium for Reliability and Reproducibility: repeated scans for test--retest studies. & T1 & 2{,}745 & 1{,}498 & SSL & -- \\
BrainLat~\cite{brainlat} & Latin American cohort with neurodegenerative diseases (AD, FTD, PD, MS) and controls. & T1, FLAIR & 632 & 587 & SSL & -- \\
M4Raw~\cite{m4raw} & Low-field (0.3\,T) multi-contrast brain MRI. & T1, T2, FLAIR & 624 & 208 & SSL & -- \\
Nigerian clinical & Clinical brain MRI acquired in Nigeria. & T1, T2, FLAIR & 283 & 82 & SSL & -- \\
\midrule
\multicolumn{7}{l}{\textbf{Used only for exclusion or not used}}\\
BHSD & Hemorrhage segmentation labels for a subset of RSNA studies; these studies are removed from pretraining. & CT & -- & -- & exclusion & -- \\
SinoCT & Reconstructed head CT series; dropped: mostly $\sim$36\,mm slabs that cannot be placed in a whole-head frame. & CT & -- & -- & excluded & -- \\
\end{longtable}
\endgroup

\subsection{Splits and hold-outs}

\begin{itemize}
  \item \textbf{Subject-level splits.} Volumes are assigned train/validation/test (75/10/15) by subject, stratified
  by modality and dataset family, so all scans of a subject share one split. Splits are inherited across store
  rebuilds so results stay comparable; official test sets are pinned where they exist (e.g.\ AISD).
  \item \textbf{Held-out families.} CQ500/Seg-CQ500, INSTANCE, AISD, ISLES 2022/2024, ATLAS, WMH, BrainMetShare and MSD
  are never used for pretraining, nor are FOMO families whose source mapping overlaps them, nor the BHSD studies.
  \item \textbf{One pretraining volume per subject and contrast.} Duplicates (e.g.\ the same BraTS subject in two
  releases, repeat OASIS sessions) are reduced to a single volume for pretraining.
  \item \textbf{Pretraining overlap is recorded.} Several downstream families (OASIS, IXI, ABIDE, AOMIC, Pixar,
  BraTS, RSNA) are part of pretraining as unlabeled images; their test subjects are never trained on, and every
  evaluation record stores the fraction of the task's volumes that were in the pretraining set.
\end{itemize}

\subsection{Downstream tasks}

\begin{table}[h]
\centering\footnotesize
\begin{tabularx}{\textwidth}{@{}L{3.0cm} l L{2.6cm} X l@{}}
\toprule
Task & Type & Data & Target & Headline metric \\
\midrule
\code{cq500} & multi-label & CQ500, Seg-CQ500 & ICH, IPH, IVH, SDH, SAH, calvarial fracture, mass effect, midline shift (reader majority) & macro AUC \\
\code{rsna\_ich} & multi-label & RSNA & any hemorrhage, epidural, intraparenchymal, intraventricular, subarachnoid, subdural (study level) & macro AUC \\
\code{oasis1\_cdr}, \code{oasis2\_cdr} & binary & OASIS-1/2 & CDR $>0$ & AUC \\
\code{oasis2\_conversion} & binary & OASIS-2 & CDR 0 at first visit, CDR $>0$ later & AUC \\
\code{abide\_asd} & binary & ABIDE I & autism spectrum disorder vs.\ control & AUC \\
\code{msd} & binary & MSD & high enhancing-tumor burden & AUC \\
\code{*\_age} (6 tasks) & regression & OASIS-1, IXI, Pixar, PIOP1/2, ID1000 & chronological age & $R^2$ \\
\midrule
\code{seg\_cq500}, \code{instance} & segmentation & Seg-CQ500, INSTANCE & intracranial hemorrhage & macro Dice \\
\code{brats2023\_gli}, \code{brats2023\_ped} & segmentation & BraTS & enhancing tumor, tumor core, whole tumor & macro Dice \\
\code{brats2023\_men} & segmentation & BraTS MEN & meningioma & macro Dice \\
\code{brats2024\_gli} & segmentation & BraTS 2024 & enhancing tumor, resection cavity & macro Dice \\
\code{atlas}, \code{isles22}, \code{isles24}, \code{aisd} & segmentation & stroke sets & ischemic lesion / infarct & macro Dice \\
\code{wmh} & segmentation & WMH & white-matter hyperintensities & macro Dice \\
\midrule
\code{probe3d} & token probes & held-out lesion sets, all families & lesion detectability (AP lift, AUROC), modality/dataset fingerprint, position $R^2$ & per probe \\
\bottomrule
\end{tabularx}
\caption{The downstream suite: 25 entries, run for every checkpoint.}
\label{tab:tasks}
\end{table}

%=====================================================================================================================
\section{Preprocessing}

Every scan is mapped to one stored representation: a $256\times256\times160$ grid of 1\,mm voxels in a common
anatomical frame, in native intensity units, with a brain mask and a record of how it got there. The steps, and why
each is there:

\begin{enumerate}
  \item \textbf{Read and orient.} Every source (NIfTI, Analyze, DICOM series) is read into one orientation convention;
  CT values are repaired where a converter shifted them and clamped to the valid Hounsfield range.
  \textit{Why:} about 30 sources disagree on file formats, axis orders and padding values; the model must not have to
  learn them.
  \item \textbf{Bias-field correction (raw MRI only).} N4~\cite{n4} removes the smooth intensity inhomogeneity of
  receive coils; releases that are already corrected are left alone. \textit{Why:} coil shading is a scanner
  property, not anatomy, and would otherwise be a shortcut and a confound for intensity statistics.
  \item \textbf{Registration to a template.} Each scan is centered, then fitted rigidly and affinely to the MNI152
  template~\cite{mni} by mutual information, which needs no intensity correspondence, so CT and every MRI contrast map
  to the same T1 template. Skull-stripped releases are fitted to the template's brain only. The fit runs on the CPU
  (SimpleITK) or on the GPU (a validated PyTorch port, about $10\times$ faster), and the full transform is stored.
  \textit{Why:} the same anatomical point should land at the same position in every scan, so the encoder's position
  encoding means anatomy and datasets are comparable. The fit is deliberately affine, not deformable, so individual
  anatomy (atrophy, ventricle size, lesions) is preserved; the stored transform recovers true physical sizes and maps
  outputs back onto the original scan.
  \item \textbf{One resampling onto a fixed grid.} The scan is resampled once, through the fitted transform, onto
  $256\times256\times160$\,mm at 1\,mm, with anti-aliasing where the source is finer. \textit{Why:} a single
  interpolation avoids compounding blur; 160\,mm of depth holds the whole brain (vertex to cerebellum, 154\,mm in
  MNI space) and is divisible by every patch depth used.
  \item \textbf{Brain mask.} SynthStrip~\cite{synthstrip}, a contrast- and modality-agnostic skull-stripping network,
  runs on the GPU; each mask is checked against the registered template brain (Dice $\geq 0.85$ within the scan's
  field of view) with a classical fallback. \textit{Why:} the mask provides parenchyma statistics for MRI intensity
  normalization, enables the skull-dropout augmentation that makes skull-on and skull-stripped scans look alike, and
  is a quality signal for the whole pipeline.
  \item \textbf{Native intensity units.} CT stays in Hounsfield units; MRI becomes a robust z-score (median and MAD
  inside the brain). Windowing is not baked in; random windows are applied during training. \textit{Why:} CT units
  are physical and should be kept; MRI units are arbitrary and need per-scan normalization; deferring windows lets one
  stored volume be seen under many renderings.
  \item \textbf{Storage and index.} One compressed file per volume holds the image and its brain mask (one file read
  per training sample); label masks are stored alongside for downstream use. An index records subjects, splits,
  hold-outs, acquisition geometry (including slice thickness) and quality-control fields. \textit{Why:} fast
  loading for training, and every downstream number can be traced to the exact data version.
  \item \textbf{Quality control and exclusions.} Volumes that fail centering or foreground checks are not stored;
  SinoCT is dropped; RSNA volumes that fail the brain-mask check (mostly mis-assembled DICOM series) are excluded from
  pretraining but kept for evaluation. Raw data is kept as compressed archives (JPEG-LS for imaging,
  byte-exact for everything else) so the pipeline can be rerun.
\end{enumerate}

%=====================================================================================================================
\section{Methods}

\subsection{Encoder}

\begin{itemize}
  \item \textbf{Architecture.} A 3D Vision Transformer (pre-norm, 12 blocks, query/key normalization), in sizes
  \emph{tiny} (width 192, 3 heads) and \emph{small} (width 384, 6 heads); \emph{base} (768, 12) is available.
  \item \textbf{Tokens.} Anisotropic patches of $16\times16\times8$\,mm; a full volume is a $20\times16\times16$ grid
  of 5{,}120 tokens. Patch size is variable (the patch embedding is resized, FlexiViT-style~\cite{flexivit}) and a
  small network adds a function of the log patch size in mm, so the model knows the physical scale of each token.
  \item \textbf{Position.} Rotary position embedding~\cite{rope} on the token centers in millimetres (3D, base
  100\,mm), so crops of any size and position share one coordinate system.
  \item \textbf{Air tokens} (patch center below a CT/MRI threshold) are dropped before the transformer; rows of
  different lengths are bucketed to limit padding.
  \item \textbf{Optional prefix tokens:} a CLS token and register tokens~\cite{registers}, without position and never
  dropped. With a CLS token, the pooled embedding is its output; otherwise it is the mean over patch tokens.
  \item \textbf{Heads used only in training:} a projector MLP for the pooled embedding, a token head (128-d) for the
  dense terms, and a small predictor transformer for masked prediction.
\end{itemize}

\subsection{Views and augmentation}

Crops are cut on the volume's own token grid, so overlapping crops share tokens exactly: two global crops
($128\times128\times64$\,mm, the second inside the first), two local crops ($64\times64\times32$\,mm) and two fine crops
($32\times32\times16$\,mm at $8\times8\times4$\,mm patches) inside a global crop, plus one crop seen under two
intensity windows. Crops must contain at least 15\% foreground. Appearance is re-rendered independently per crop, so
that views can only be matched through anatomy, never through contrast, resolution or the presence of a skull:
random intensity windows; a random monotone contrast curve with knots at tissue quantiles, inverted with probability
0.3 (as T1 and T2 order tissues differently); a smooth MRI bias field; blur along the acquisition axis up to a
random slice thickness (using the recorded native thickness); isotropic blur, slow drift and noise; and, with
probability 0.5, \emph{skull dropout}, cutting a whole volume to its brain mask (only masks that passed quality
control).

\subsection{Objective}

With $\lambda=0.02$,
\begin{equation*}
\mathcal{L} = (1-\lambda)\,\mathcal{L}_{\text{inv}} + \lambda\,\mathcal{L}_{\text{SIGReg}} + w_{\text{win}}\mathcal{L}_{\text{win}}
+ w_{\text{dense}}\big[(1-\lambda)\mathcal{L}_{\text{dense}} + \lambda\,\mathcal{L}_{\text{SIGReg}}^{\text{tok}}\big]
+ w_{\text{pred}}\big(\mathcal{L}_{\text{pred}} + w_{\text{pix}}\mathcal{L}_{\text{pix}}\big).
\end{equation*}
\begin{itemize}
  \item $\mathcal{L}_{\text{inv}}$, $\mathcal{L}_{\text{SIGReg}}$: the LeJEPA terms on pooled embeddings: every view is
  pulled to the mean of the global views, and the batch of embeddings is pushed toward an isotropic Gaussian with a
  sketched Epps--Pulley test on random 1-D projections. $\mathcal{L}_{\text{win}}$ ties the pair of windows.
  \item $\mathcal{L}_{\text{dense}}$: every grid position covered by two or more views is pulled to the mean of its
  token embeddings across views (fine tokens averaged into their coarse position);
  $\mathcal{L}_{\text{SIGReg}}^{\text{tok}}$ applies SIGReg to 32 sampled token embeddings per view.
  \item $\mathcal{L}_{\text{pred}}$, $\mathcal{L}_{\text{pix}}$: three blocks of 3--6 tokens are hidden in each global
  crop; a 4-layer predictor regresses their token-head embeddings from the unmasked crop (stop-gradient target) and
  the $2\times2\times2$ cell means of their voxels. The weight ramps up over 2{,}000 steps.
\end{itemize}

\subsection{Training}

The \emph{small} recipe: AdamW (lr $2.5\times10^{-4}$, $\beta=(0.9, 0.95)$, weight decay 0.05, none on norms, biases
and prefix tokens), 2{,}000 warm-up steps then cosine decay, gradient clipping at 1.0, bfloat16, 30{,}000 steps,
batch 96 volumes per GPU on 2 GPUs. Sampling is balanced across dataset families with mass $\propto n^{0.5}$, so the
largest families (FOMO Yale, RSNA) do not dominate. Checkpoints record their own provenance (code version, store
format and grid, index hash) so evaluations can name exactly what a model was trained on.

\subsection{Frozen evaluation protocol}

\begin{itemize}
  \item \textbf{Embedding.} Each volume is encoded once, whole, under the evaluation windows (CT: brain, subdural,
  bone; MRI: tissue, wide), giving a 5{,}120-token map, a foreground mask and, for CLS models, the CLS embedding.
  \item \textbf{Classification and regression heads.} PCA (16 components, fitted on training tokens) followed by
  three readouts: per-axial-slab mean and max, a $4\times4\times4$ grid of cell means, and global mean and max, each
  with logistic or ridge regression; a small 3D convolutional probe on the full-width token map (seed ensemble,
  early stopping); and a linear head on the CLS embedding when present. Regularization is chosen on validation.
  \item \textbf{Segmentation.} A token U-Net decoder trained on cached token maps upsamples to one logit per voxel
  (cross-entropy plus soft Dice per region); training volumes are also encoded shifted off the patch grid and flipped.
  Dice is reported on the 64 axial slices scored by the 2D reference and on the whole volume.
  \item \textbf{Token probes} (\code{probe3d}). Lesion detection from single tokens (linear and MLP, average-precision
  lift over prevalence and AUROC, per held-out lesion family); fingerprinting (balanced accuracy of predicting
  modality and source dataset from pooled embeddings; lower is better for dataset); and location (ridge regression from
  a token to its normalized position, $R^2$, plus the participation ratio of the token covariance). Pixel-statistics
  and random-initialization controls are run alongside.
  \item \textbf{Statistics and selection.} Bootstrap 95\% intervals over test volumes; for each model and task the
  head is chosen on validation, never on test. Results are compared across models with a suite-wide score: each
  task's headline metric expressed as gain over the random-initialization floor, averaged within task groups (CT
  classification, MRI classification, regression, segmentation, probes) and then across groups, so the many
  segmentation tasks do not dominate.
\end{itemize}

\subsection{Task-conditioned post-training}

\begin{itemize}
  \item \textbf{Conditioning.} A learned task vector enters the encoder as an extra prefix token (the same mechanism as
  CLS and registers); FiLM-style modulation is a later alternative. Each task keeps its own small output head
  (classification, regression or segmentation decoder).
  \item \textbf{Task granularity.} One vector per label (each hemorrhage subtype, age, CDR) with a shared
  per-dataset component, so related labels can share structure; batches mix tasks.
  \item \textbf{Null task.} The task vector is replaced by a learned null vector for 10--20\% of samples, so the model
  keeps a task-free embedding that can be compared like-for-like with the self-supervised model.
  \item \textbf{Balancing.} Temperature-based task sampling and per-task loss normalization, so RSNA (about 25{,}000
  labeled volumes) does not drown OASIS (about 400).
  \item \textbf{Preserving generality.} The self-supervised loss is kept as an auxiliary term (or only adapters/LoRA
  are trained), and the token probes are rerun before and after post-training.
\end{itemize}

%=====================================================================================================================
\section{Planned experiments}

All experiments use the v6 store, the downstream suite of Table~\ref{tab:tasks}, two seeds where differences are
small, and the suite-wide score. A change is adopted when it improves the suite score beyond the seed spread and no
task group regresses beyond its confidence interval.

\begin{table}[h]
\centering\footnotesize
\begin{tabular}{@{}L{0.6cm} L{3.4cm} L{5.4cm} L{5.6cm}@{}}
\toprule
& Question & Arms & Readout \\
\midrule
E0 & Floor and baseline & random-init encoder; current \emph{small} recipe & reference rows for every comparison \\
E1 & Prefix tokens & none (mean pooling) / 1 CLS / CLS + 4 registers & global tasks (CLS vs.\ token heads), dense tasks and probes (registers) \\
E2 & Token-level SIGReg & 4 vs.\ 32 sampled tokens per view (32 adopted) & token-probe quality, segmentation \\
E3 & Full-volume context & + one full-volume view per sample & global tasks; train/test context gap \\
E4 & Mixed patch size & per-crop patch $16\times16\times8$ or $8\times8\times4$, with cross-scale view pairs & dense tasks at both scales; coarse-token quality \\
E5 & Scale, length and data & tiny / small / base; longer schedules; sampling exponent; FOMO share; QC exclusions & suite score vs.\ compute \\
E6 & Task-conditioned post-training & see Table~\ref{tab:ladder} & seen and held-out tasks; label efficiency; probes before/after \\
\bottomrule
\end{tabular}
\caption{Experiment plan. E0 and E1 are queued on the v6 store; E2's setting is already adopted.}
\label{tab:exp}
\end{table}

\paragraph{Post-training comparison (E6).} Roughly a third of the tasks are fixed as \emph{held out of
post-training} before any post-training run (rotated if compute allows). On seen tasks, the rungs of
Table~\ref{tab:ladder} separate what the representation, head capacity, supervision and conditioning each
contribute. On held-out tasks, the self-supervised and post-trained backbones are compared with identical heads, and
a new task vector is learned with everything else frozen (a few hundred parameters). Heads are also trained on 10\%
and 100\% of each task's labels.

\begin{table}[h]
\centering\footnotesize
\begin{tabular}{@{}l L{7.5cm} L{5.5cm}@{}}
\toprule
Rung & Model & Isolates \\
\midrule
1 & frozen self-supervised encoder + linear head & the representation \\
2 & frozen self-supervised encoder + small attentive / MLP head & head capacity \\
3 & single-task fine-tuning, one model per task & supervised ceiling per task \\
4 & multi-task fine-tuning, no task vector & shared supervision \\
5 & multi-task fine-tuning with task vector & conditioning \\
\bottomrule
\end{tabular}
\caption{Comparison ladder for task-conditioned post-training, on identical splits.}
\label{tab:ladder}
\end{table}

%=====================================================================================================================
\section{Known limitations}

\begin{itemize}
  \item \textbf{Affine size normalization.} Stored lengths and volumes are scaled by the registration (volumes by a
  median factor of 1.35 in our data). True sizes are recovered from the stored transform; size-sensitive endpoints
  must use it or condition on it.
  \item \textbf{RSNA DICOM assembly.} About 2\% of RSNA studies mix slice thicknesses or repeat positions within a
  series, and gantry tilt is only removed when the affine fit succeeds. Affected volumes are excluded from pretraining;
  the reader will be replaced (position-based assembly with tilt correction) before the next rebuild.
  \item \textbf{Mixed registration engines.} Volumes are registered by SimpleITK or by the GPU port depending on the
  processing phase; each volume records which, and the GPU port was validated as equal or better except on heavily
  tilted CT.
  \item \textbf{Shared pretraining and evaluation images.} Several downstream families are in the unlabeled
  pretraining pool (Section~3.2); held-out families provide the clean test of transfer.
  \item \textbf{OASIS-1.} SynthStrip under-segments the atlas-registered OASIS-1 images; their masks fail quality
  control and are not used for skull dropout.
\end{itemize}

%=====================================================================================================================
\begin{thebibliography}{99}\small
\bibitem{lejepa} R.~Balestriero and Y.~LeCun. LeJEPA: provable and scalable self-supervised learning without the heuristics. \emph{arXiv}, 2025.
\bibitem{flexivit} L.~Beyer et al. FlexiViT: one model for all patch sizes. \emph{CVPR}, 2023.
\bibitem{rope} J.~Su et al. RoFormer: enhanced transformer with rotary position embedding. \emph{Neurocomputing}, 2024.
\bibitem{registers} T.~Darcet, M.~Oquab, J.~Mairal and P.~Bojanowski. Vision transformers need registers. \emph{ICLR}, 2024.
\bibitem{n4} N.~J.~Tustison et al. N4ITK: improved N3 bias correction. \emph{IEEE Trans.\ Med.\ Imaging}, 2010.
\bibitem{mni} V.~Fonov et al. Unbiased average age-appropriate atlases for pediatric studies. \emph{NeuroImage}, 2011.
\bibitem{synthstrip} A.~Hoopes et al. SynthStrip: skull-stripping for any brain image. \emph{NeuroImage}, 2022.
\bibitem{rsna} A.~E.~Flanders et al. Construction of a machine learning dataset through collaboration: the RSNA 2019 brain CT hemorrhage challenge. \emph{Radiology: Artificial Intelligence}, 2020.
\bibitem{cq500} S.~Chilamkurthy et al. Deep learning algorithms for detection of critical findings in head CT scans: a retrospective study. \emph{The Lancet}, 2018.
\bibitem{aisd} K.~Liang et al. Symmetry-enhanced attention network for acute ischemic infarct segmentation with non-contrast CT images. \emph{MICCAI}, 2021.
\bibitem{brats21} U.~Baid et al. The RSNA-ASNR-MICCAI BraTS 2021 benchmark on brain tumor segmentation and radiogenomic classification. \emph{arXiv}, 2021.
\bibitem{bratsmen} D.~LaBella et al. The ASNR-MICCAI brain tumor segmentation (BraTS) challenge 2023: intracranial meningioma. \emph{arXiv}, 2023.
\bibitem{bratsped} A.~F.~Kazerooni et al. The brain tumor segmentation (BraTS) challenge 2023: focus on pediatrics. \emph{arXiv}, 2023.
\bibitem{msd} M.~Antonelli et al. The Medical Segmentation Decathlon. \emph{Nature Communications}, 2022.
\bibitem{brainmetshare} E.~Gr{\o}vik et al. Deep learning enables automatic detection and segmentation of brain metastases on multisequence MRI. \emph{J.\ Magn.\ Reson.\ Imaging}, 2020.
\bibitem{atlas} S.-L.~Liew et al. A large, curated, open-source stroke neuroimaging dataset to improve lesion segmentation algorithms. \emph{Scientific Data}, 2022.
\bibitem{isles22} M.~R.~Hernandez Petzsche et al. ISLES 2022: a multi-center magnetic resonance imaging stroke lesion segmentation dataset. \emph{Scientific Data}, 2022.
\bibitem{wmh} H.~J.~Kuijf et al. Standardized assessment of automatic segmentation of white matter hyperintensities and results of the WMH segmentation challenge. \emph{IEEE Trans.\ Med.\ Imaging}, 2019.
\bibitem{oasis1} D.~S.~Marcus et al. Open Access Series of Imaging Studies (OASIS): cross-sectional MRI data in young, middle aged, nondemented, and demented older adults. \emph{J.\ Cogn.\ Neurosci.}, 2007.
\bibitem{oasis2} D.~S.~Marcus et al. Open Access Series of Imaging Studies: longitudinal MRI data in nondemented and demented older adults. \emph{J.\ Cogn.\ Neurosci.}, 2010.
\bibitem{abide} A.~Di~Martino et al. The autism brain imaging data exchange: towards a large-scale evaluation of the intrinsic brain architecture in autism. \emph{Molecular Psychiatry}, 2014.
\bibitem{aomic} L.~Snoek et al. The Amsterdam Open MRI Collection, a set of multimodal MRI datasets for individual difference analyses. \emph{Scientific Data}, 2021.
\bibitem{pixar} H.~Richardson et al. Development of the social brain from age three to twelve years. \emph{Nature Communications}, 2018.
\bibitem{hbn} L.~M.~Alexander et al. An open resource for transdiagnostic research in pediatric mental health and learning disorders. \emph{Scientific Data}, 2017.
\bibitem{nki} K.~B.~Nooner et al. The NKI-Rockland sample: a model for accelerating the pace of discovery science in psychiatry. \emph{Frontiers in Neuroscience}, 2012.
\bibitem{corr} X.-N.~Zuo et al. An open science resource for establishing reliability and reproducibility in functional connectomics. \emph{Scientific Data}, 2014.
\bibitem{brainlat} P.~Prado et al. The BrainLat project, a multimodal neuroimaging dataset of neurodegeneration from underrepresented backgrounds. \emph{Scientific Data}, 2023.
\bibitem{m4raw} M.~Lyu et al. M4Raw: a multi-contrast, multi-repetition, multi-channel MRI k-space dataset for low-field MRI research. \emph{Scientific Data}, 2023.
\end{thebibliography}

\end{document}
